% Auto-generated targeting coverage analysis table
\begin{table}[H]
  \centering
  \scriptsize
  \setlength{\tabcolsep}{2pt}
  \renewcommand{\arraystretch}{1.1}
  \begin{tabular}{lcccccccc}
    \toprule
    DGP & N=25 & N=50 & N=100 & N=200 & N=400 & N=800 & N=1600 & N=3200 \\
    \midrule
    TN & (93.3, 94.6) & (94.4, 96.0) & (94.9, 95.2) & (95.8, 95.4) & (94.7, 94.7) & (93.9, 95.1) & (95.3, 95.2) & (94.3, 94.3) \\
    GS3 & (95.1, 95.9) & (95.2, 95.8) & (95.7, 95.9) & (95.5, 94.9) & (94.3, 95.1) & (95.3, 94.9) & (95.1, 94.6) & (96.3, 96.1) \\
    GA3 & (94.8, 95.4) & (94.7, 95.3) & (95.6, 94.6) & (95.3, 95.0) & (94.9, 95.2) & (94.3, 94.2) & (96.2, 94.9) & (94.8, 94.6) \\
    GS5 & (95.5, 96.0) & (94.3, 94.7) & (96.2, 96.0) & (94.7, 94.5) & (93.1, 93.6) & (94.9, 94.8) & (95.5, 94.4) & (94.6, 94.0) \\
    Step & (93.5, 95.1) & (95.6, 96.4) & (92.6, 93.7) & (95.6, 95.3) & (94.0, 94.6) & (94.0, 95.1) & (95.2, 95.0) & (94.4, 94.3) \\
    Sine & (92.8, 93.6) & (95.1, 95.1) & (94.5, 94.7) & (95.8, 95.4) & (94.7, 94.9) & (94.3, 95.3) & (95.3, 95.3) & (94.4, 94.1) \\
    \bottomrule
  \end{tabular}
  \caption{Coverage for mean}
  \label{tab:coverage_targeting_mean}
\end{table}